\documentclass[11pt]{article}
\usepackage[margin=1in]{geometry}
\usepackage{amsmath,amssymb,amsthm,graphicx,microtype}
\usepackage[colorlinks=true,linkcolor=blue,urlcolor=blue,citecolor=blue]{hyperref}
\newtheorem{theorem}{Theorem}
\newtheorem{lemma}{Lemma}
\newcommand{\R}{\mathbb R}
\newcommand{\cQ}{Q}
\newcommand{\Ps}{P_s}
\title{Square anisotropy: the Wulff square is never fractional-isoperimetric\\
\large A sharp partial result for the open problem in arXiv:1304.0699}
\author{Autonomous proof audit --- lane 9}
\date{17 August 2026}

\begin{document}
\maketitle

\begin{center}
\fbox{\parbox{0.91\linewidth}{\textbf{Candidate partial result.}
For every $0<s<1$, the square unit ball is not a minimizer of its own
anisotropic fractional perimeter: its nonlocal mean curvature varies strictly
along each side.  The conclusion transfers affinely to every parallelogram.
In contrast, the square is the unique minimizer among all fixed-area aligned
rectangles.  Thus the true optimizer must leave the most obvious Wulff family.}}
\end{center}

\section{Source problem and current frontier}

For an origin-symmetric convex body $K\subset\R^n$ and $0<s<1$, Ludwig
defines
\[
 \Ps(E,K)=\int_E\int_{E^c}\frac{dx\,dy}{\lVert x-y\rVert_K^{n+s}}.
\]
The anisotropic fractional isoperimetric problem is to determine the equality
cases in
\[
 \Ps(E,K)\ge\gamma_s(K)|E|^{(n-s)/n}.
\]
The source explicitly says that determination of the minimizers remains open.

\begin{figure}[ht]
\centering
\includegraphics[width=0.96\linewidth]{figures/open_problem_crop.png}
\caption{The open minimizer problem in arXiv:1304.0699, printed page 4.}
\end{figure}

The Euclidean case is solved by balls, and an invertible linear change of
variables solves every ellipsoidal anisotropy.  For nonellipsoidal $K$, later
work proves structural results but not a classification: Kreuml proves an
unconditional star-body conclusion under strict convexity and
unconditionality \cite{Kreuml}, while Cai derives a necessary fractional-area
measure condition for convex optimizers \cite{Cai}.  The result below gives an
explicit all-$s$ obstruction for the most elementary crystalline anisotropy.

\section{The Wulff square is not stationary}

Let
\[
 \cQ=[-1,1]^2,
 \qquad \lVert z\rVert_{\cQ}=\max\{|z_1|,|z_2|\},
 \qquad k_s(z)=\lVert z\rVert_{\cQ}^{-2-s}.
\]
At the regular boundary point $p_t=(1,t)$, $|t|<1$, define the principal-value
nonlocal curvature
\[
 H_s(t)=\operatorname{PV}\int_{\R^2}
  \bigl(\mathbf1_{\cQ^c}(y)-\mathbf1_{\cQ}(y)\bigr)
  k_s(y-p_t)\,dy.
\]

\begin{lemma}\label{lem:curvature}
The function $H_s$ is even and is strictly increasing on $(0,1)$.  More
precisely,
\begin{equation}\label{eq:derivative}
 H_s'(t)=2\int_{-2}^{0}
 \left[
  \max\{|x|,1-t\}^{-2-s}
  -\max\{|x|,1+t\}^{-2-s}
 \right]dx>0.
\end{equation}
\end{lemma}

\begin{proof}
After the change of variables $z=y-p_t$, the translated interior is
\[
 R_t=\cQ-p_t=[-2,0]\times[-1-t,1-t].
\]
In the difference $H_s(t)-H_s(0)$ the singular half-space contributions
cancel, and hence
\[
 H_s(t)-H_s(0)
 =-2\int_{\R^2}
   (\mathbf1_{R_t}-\mathbf1_{R_0})(z)k_s(z)\,dz.
\]
The symmetric difference stays away from the origin to first order, so this
finite expression may be differentiated by moving the two horizontal edges.
If
\[
 J(t)=\int_{-2}^{0}\int_{-1-t}^{1-t}k_s(x,y)\,dy\,dx
\]
is interpreted only through differences, then
\[
 J'(t)=\int_{-2}^{0}
 \bigl(k_s(x,1+t)-k_s(x,1-t)\bigr)\,dx.
\]
Since $H_s'=-2J'$, this is \eqref{eq:derivative}.  The integrand there is
nonnegative and is strictly positive, for example, whenever $|x|<1-t$.
Reflection in the horizontal axis gives evenness.
\end{proof}

\begin{theorem}\label{thm:notminimal}
For every $0<s<1$, no translate or dilate of $\cQ$ minimizes
$\Ps(\,\cdot\,,\cQ)$ under a fixed-area constraint.  More generally, if $K$
is any origin-symmetric parallelogram, no homothetic copy of $K$ is a
minimizer of $\Ps(\,\cdot\,,K)$.
\end{theorem}

\begin{proof}
A volume-constrained minimizer is stationary under smooth normal variations
supported in the relative interior of its sides; its nonlocal curvature must
therefore equal one Lagrange multiplier there.  Here this necessary condition
can also be used constructively.  Choose two small side intervals, one near
$t=0$ and one near a fixed $t_0\in(0,1)$, and a smooth function $\varphi$
supported on them with zero integral, positive near $0$, and negative near
$t_0$, with equal masses.  Replace the right boundary $x=1$ by
$x=1+\varepsilon\varphi(t)$.  The area is unchanged exactly, and the standard
first-variation formula gives
\[
 \left.\frac d{d\varepsilon}\right|_{\varepsilon=0}
 \Ps(E_\varepsilon,\cQ)
 =\int_{-1}^{1}H_s(t)\varphi(t)\,dt<0
\]
by Lemma~\ref{lem:curvature}.  Thus $\cQ$ is not even stationary.

If $A$ is invertible, then
\[
 \lVert A(x-y)\rVert_{A\cQ}=\lVert x-y\rVert_{\cQ},
 \quad
 \Ps(AE,A\cQ)=|\det A|^2\Ps(E,\cQ),
 \quad |AE|=|\det A|\,|E|.
\]
Every origin-symmetric parallelogram is $A\cQ$ for some $A$, so the
nonminimality transfers.  Translation and dilation invariance give the stated
homothetic version.
\end{proof}

\section{Nevertheless, the square is best among rectangles}

The previous theorem is not a mere aspect-ratio instability.  Let
\[
 R_q=[-q/2,q/2]\times[-1/(2q),1/(2q)],\qquad q\ge1,
\]
so every $R_q$ has area one.

\begin{theorem}\label{thm:rectangles}
For every $0<s<1$, the map
$q\mapsto\Ps(R_q,\cQ)$ is strictly increasing on $(1,\infty)$.  Hence the
square $R_1$ is the unique minimizer, up to interchanging the axes, among all
area-one axis-aligned rectangles.
\end{theorem}

\begin{proof}
For a rectangle with side lengths $A\ge B$, its covariogram is
\[
 g_{A,B}(z)=(A-|z_1|)_+(B-|z_2|)_+.
\]
The translation formula and the layer-cake identity
$r^{-2-s}=(2+s)\int_r^\infty t^{-3-s}\,dt$ give
\begin{equation}\label{eq:layer}
 \Ps(R,\cQ)=(2+s)\int_0^\infty t^{-3-s}
 \bigl(4ABt^2-h_A(t)h_B(t)\bigr)\,dt,
\end{equation}
where
\[
 h_C(t)=
 \begin{cases}2Ct-t^2,&0\le t\le C,\\ C^2,&t\ge C.\end{cases}
\]
Indeed, the bracket in \eqref{eq:layer} is
$\int_{[-t,t]^2}(|R|-g_{A,B}(z))\,dz$.

Substitute $A=q$, $B=q^{-1}$, split the integral at $B$ and $A$, and integrate
the resulting powers.  Direct simplification yields
\begin{align*}
 \frac d{dq}\Ps(R_q,\cQ)
  &=\frac{2(s+2)q^{-s-3}}{s(1-s)(1+s)}G_s(q),\\
 G_s(q)&=2s q^{2+2s}-(s+1)q^{2s}+1-s.
\end{align*}
Now $G_s(1)=0$ and
\[
 G_s'(q)=2s(s+1)q^{2s-1}(2q^2-1)>0
 \qquad(q\ge1).
\]
The claimed strict increase follows.
\end{proof}

Theorem~\ref{thm:rectangles} and Theorem~\ref{thm:notminimal} together force
any true optimizer for square anisotropy to be a genuinely nonrectangular
shape.  Affine covariance gives the analogous statement for parallelogram
anisotropy and aligned parallelogram competitors.

\section{Upgrade audit, scope, and novelty}

Several routes toward a full classification were tested.  Ellipsoids reduce
to the Euclidean theorem, but anisotropic Riesz rearrangement is unavailable
outside that case.  Steiner symmetrization reaches unconditional star bodies
only under additional hypotheses.  The $s\uparrow1$ limit selects the moment
body but does not identify any fixed-$s$ optimizer.  Finally,
Lemma~\ref{lem:curvature} turns the square problem into a nonlinear constant
nonlocal-curvature equation for an unknown rounded box body; neither the
rectangle calculation nor current Alexandrov-type theorems for radial kernels
control that anisotropic equation.  These are genuine obstructions, so no full
classification is claimed.

Bounded searches using the exact source phrase, title, square/cube,
parallelogram, Wulff shape, optimizer, and anisotropic nonlocal curvature found
the source, Kreuml's structural theorem, and Cai's recent necessary condition,
but no all-$s$ square nonstationarity calculation or exact rectangle comparison.
Novelty is therefore plausible but unverified.  Expert review should focus on
the first-variation step at a polygonal boundary and the layer-cake constants.

\begin{thebibliography}{9}
\bibitem{Ludwig}
M. Ludwig,
\emph{Anisotropic fractional perimeters},
J. Differential Geom. \textbf{96} (2014), 77--93;
arXiv:1304.0699.

\bibitem{Kreuml}
A. Kreuml,
\emph{The anisotropic fractional isoperimetric problem with respect to
unconditional unit balls},
Commun. Pure Appl. Anal. \textbf{20} (2021), 783--799,
\href{https://doi.org/10.3934/cpaa.2020290}{doi:10.3934/cpaa.2020290}.

\bibitem{Cai}
X. Cai,
\emph{Anisotropic fractional area measures},
arXiv:2510.05279 (2025), Theorem 6.1.
\end{thebibliography}

\end{document}
